% The logical content: what the checker's guards mean, stated as propositions.
%
% REQUIRED PREAMBLE (declare in main.tex if not already present):
%   \usepackage{amsmath, amssymb, amsthm}
%   \newtheorem{proposition}{Proposition}
%   \newtheorem{corollary}[proposition]{Corollary}
%
% Every proposition below was machine-checked against z3 before being written;
% the script is in the artefact. Proofs are short by design -- the content is
% that these are the properties the system's guards exist to preserve, not that
% any one of them is difficult.

\section{Formalisation \& logic}
\label{sec:logic}

The rule language is quantifier-free first-order logic over linear arithmetic:
propositional connectives whose atoms are arithmetic comparisons on the reals
and integers, rather than bare propositions. It is the SMT fragment
$\mathrm{QF\_LIRA}$, and is decidable. Two constructs are added --- a bounded
past-time temporal operator, flattened to derived state, and a relative-change
operator over a value at an earlier event --- and two things are deliberately
absent: quantifiers, and arithmetic relating two variables.

No rule contains a quantifier; every variable is free, which is what allows a
satisfying assignment to \emph{be} a patient. The quantification sits one level
up, in the verdict: a patient record is a \emph{partial} interpretation, and the
two definite verdicts are universal statements over the interpretations
extending it --- every completion violates, or none does --- each discharged as a
single unsatisfiability check. That is the whole reason a third verdict is
needed (Propositions~\ref{prop:determinacy} and~\ref{prop:onequery}).

Deontic force is a dyadic annotation selecting a rule's violation condition, not
a modal operator over formulas; it is never nested, which keeps the standard
deontic paradoxes out of scope. Each proposition below was machine-checked
before being stated.

\paragraph{Setup.}
A rule $R$ denotes the material implication
$\Phi_R = \llbracket \mathit{when} \rrbracket \rightarrow
          \llbracket \mathit{then} \rrbracket$
over a signature of typed variables.

\paragraph{Derived symbols.}
Two constructs are interpreted over symbols the
signature induces rather than ones a rule names directly. For each declared
clinical event $e$: $\mathsf{occ}_e$, a boolean recording that $e$ \textbf{occ}urred, and
$\mathsf{dse}_e$, an integer counting days since it did (\textbf{d}ays \textbf{s}ince \textbf{e}vent). For each event $e$ and
numeric variable $v$: $v@e$, the value $v$ held at the time of $e$. A recency
window ``$e$ within the last $n$ days'' lowers to
$\mathsf{occ}_e \wedge \mathsf{dse}_e \le n$, and a relative change ``reduce $v$
by $p$'' to $v \le (1-p)\cdot v@e$. These symbols are produced by the lowering
and never written by the generator, so each has exactly one definition.

\paragraph{The domain theory.}
The domain, $\mathcal{D}$, collects plausible physiological ranges for bound
numeric variables, code membership for coded ones, $0 \le \mathsf{dse}_e \le H$
for a fixed horizon $H$ (ten years in the implementation), and $v@e > 0$ for
each induced baseline. Unbound symbols are deliberately unconstrained --- an
axiom about a predicate the source never defined would be a claim the source
does not make. A patient record is a \emph{partial} assignment $P$: it pins the values the
record supplies and leaves every other symbol free, either because nothing
measured it or because it is an undefined predicate no data source can supply.
A \emph{completion} of the record is a model of $\mathcal{D} \wedge P$ --- one
consistent way of filling in everything the record left open. A record therefore
does not describe one patient but the set of patients compatible with it, and
the three verdicts below are statements about that whole set: every completion
violates, none does, or the record does not settle the question.
$\mathrm{Vars}(\Phi)$ denotes the symbols occurring in $\Phi$.

\[
\mathrm{verdict}(R,P) =
\begin{cases}
  \textsc{must-violate}       & \mathcal{D} \wedge P \wedge \Phi_R \ \text{unsat} \\
  \textsc{certified-compliant}& \mathcal{D} \wedge P \wedge \neg\Phi_R \ \text{unsat} \\
  \textsc{may-violate}        & \text{otherwise.}
\end{cases}
\]

\begin{proposition}[Well-posedness]
\label{prop:exhaustive}
If $\mathcal{D} \wedge P$ is satisfiable, the three above cases are mutually exclusive
and exhaustive.
\end{proposition}
\begin{proof}
Exhaustive by construction. If the first two both held, $\mathcal{D} \wedge P$
would entail both $\Phi_R$ and $\neg\Phi_R$, contradicting satisfiability.
\end{proof}

\subsection{Indeterminacy is a property of the record, not of the rule}

\begin{proposition}[Determinacy]
\label{prop:determinacy}
If $P$ assigns every symbol of $\mathrm{Vars}(\Phi_R)$ and $\mathcal{D} \wedge P$
is satisfiable, then $\mathrm{verdict}(R,P) \neq \textsc{may-violate}$.
Equivalently: \textsc{may-violate} implies some symbol of $\Phi_R$ is unpinned.
\end{proposition}
\begin{proof}
$P$ fixes a truth value for $\Phi_R$, since every symbol in $\Phi_R$ is pinned. If that value is true then $\mathcal{D}\wedge P\wedge\neg\Phi_R$ is
unsat; if false then $\mathcal{D}\wedge P\wedge\Phi_R$ is unsat. By satisfiability, one of these has to be true, so \textsc{may-violate} never occurs.
\end{proof}

The converse fails, and the failure is worth stating because it is the reason
``count the free variables'' is not a usable proxy for uncertainty. With
$\Phi = (k > 5 \wedge h) \rightarrow \neg\,\mathit{acei}$ (ACE inhibitor) and $h$ unpinned,
$k = 4$ yields \textsc{certified-compliant}: the antecedent is already false, so
$h$ is free but irrelevant. The exact criterion is not freeness but
\emph{decidedness} -- $\mathrm{verdict} = \textsc{may-violate}$ precisely when
$\mathcal{D} \wedge P$ entails neither $\Phi_R$ nor $\neg\Phi_R$ -- which is the
definition unfolded, and is what the two solver calls compute.

The bridge to the system is that an \emph{undefined predicate} is a symbol no
data source can ever pin. Where such a symbol is relevant, the rule reports
\textsc{may-violate} for every record, permanently, until a site supplies it.

\subsection{Why substituting a value for an undefined predicate is unsafe}

\begin{proposition}[Monotonicity of pinning]
\label{prop:monotone}
Let $P' \supseteq P$ pin strictly more symbols, with $\mathcal{D}\wedge P'$
satisfiable. Then
\textup{(i)} $\mathrm{verdict}(R,P)=\textsc{must-violate} \Rightarrow
\mathrm{verdict}(R,P')=\textsc{must-violate}$;
\textup{(ii)} likewise for \textsc{certified-compliant};
\textup{(iii)} $\mathrm{verdict}(R,P')=\textsc{may-violate} \Rightarrow
\mathrm{verdict}(R,P)=\textsc{may-violate}$.
\end{proposition}
\begin{proof}
Writing $\mathrm{Mod}(\chi)$ for the set of completions satisfying $\chi$,
pinning more symbols can only remove completions:
$\mathrm{Mod}(\mathcal{D}\wedge P') \subseteq \mathrm{Mod}(\mathcal{D}\wedge P)$.
An unsatisfiable conjunction stays unsatisfiable under a subset, giving (i) and
(ii); (iii) is their contrapositive.
\end{proof}

\begin{corollary}[Stubbing manufactures certainty]
\label{cor:stub}
Substituting any value for an undefined predicate can only move a verdict
\emph{out of} \textsc{may-violate}; it never produces \textsc{may-violate}, and
it never moves a verdict that was already definite.
\end{corollary}

The direction of the move is chosen by the substituted value, not by the
evidence. Table~\ref{tab:stub-machine} is machine-checked on the rule above:
the same record, the same missing datum, and two opposite definite answers.

\begin{table}[h]
  \centering
  \caption{The same undefined predicate $h$, stubbed two ways. Neither column is
  supported by the source, which states $h$ and does not define it.}
  \label{tab:stub-machine}
  \begin{tabular}{llll}
    \toprule
    Record & $h$ free & $h := \textsf{false}$ & $h := \textsf{true}$ \\
    \midrule
    $k=4.0,\ \mathit{acei}$ & \textsc{cert} & \textsc{cert} & \textsc{cert} \\
    $k=5.5,\ \mathit{acei}$ & \textsc{may}  & \textsc{cert} & \textsc{must} \\
    $k=6.0,\ \mathit{acei}$ & \textsc{may}  & \textsc{cert} & \textsc{must} \\
    \bottomrule
  \end{tabular}
\end{table}

Corollary~\ref{cor:stub} is the formal content of the ablation in
\S\ref{sec:necessity}: the measured rates say how often the move occurs on a
rule set, and the corollary says that whenever it occurs it is an assertion the
source did not make.

\subsection{Why the third verdict cannot be dropped}

\begin{proposition}[One query is not enough]
\label{prop:onequery}
Let $\mathit{viol} = \neg\Phi_R$. Then
$\mathrm{verdict} = \textsc{certified-compliant}$ iff
$\mathcal{D}\wedge P\wedge \mathit{viol}$ is unsat, while both
\textsc{must-violate} and \textsc{may-violate} make it satisfiable. Hence no
procedure whose only test is the satisfiability of
$\mathcal{D}\wedge P\wedge \mathit{viol}$ can distinguish them.
\end{proposition}
\begin{proof}
\textsc{must-violate} is $\mathcal{D}\wedge P \models \mathit{viol}$, which
implies satisfiability whenever $\mathcal{D}\wedge P$ is satisfiable;
\textsc{may-violate} asserts it directly.
\end{proof}

A two-valued checker therefore reports ``some completion of this record
violates'' in the same words as ``every completion violates'' --- the first is a
missing datum, the second is a deviation.

\subsection{Why deontic force is an operator and not a label}

Each rule carries a \emph{deontic force} $\delta \in \{\mathsf{Obl},
\mathsf{Per}\}$: an obligation ($\mathsf{Obl}$) requires the action of a patient the rule applies to, whereas a permission ($\mathsf{Per}$) offers it, so a patient who declines has deviated from nothing. The force is a conditional operator, not an annotation, and it enters the logic at exactly one point --- the condition under which the rule counts as broken:
\[
  \mathrm{viol}(R) =
  \begin{cases}
    \llbracket \mathit{when} \rrbracket \wedge
      \neg\llbracket \mathit{then} \rrbracket & \delta = \mathsf{Obl} \\
    \bot & \delta = \mathsf{Per.}
  \end{cases}
\]

\begin{proposition}[Permissions accuse no one]
\label{prop:permission}
If $\delta = \mathsf{Per}$ then $\mathcal{D}\wedge P\wedge\mathrm{viol}(R)$ is
unsatisfiable.
\end{proposition}
\begin{proof}
Immediate: the conjunction contains $\bot$.
\end{proof}

The guarantee is structural rather than a condition applied to a verdict string
afterwards --- there is no record and no completion of its unknowns for which a
permission yields a violation. Replacing $\bot$ by the obligation's condition,
as any lowering that reads a single \texttt{Violated} expression must, moves the
verdict for four permissive rule--record pairs in five (\S\ref{sec:necessity}).

\subsection{Two lowerings that are wrong without a side condition}

\begin{proposition}[The occurrence flag is not redundant]
\label{prop:occ}
$\mathcal{D}$ contains $0 \le \mathsf{dse}_e \le H$. Under any interpretation
with $\mathsf{occ}_e = \textsf{false}$, the formula $\mathsf{dse}_e \le n$ is
satisfiable, while $\mathsf{occ}_e \wedge \mathsf{dse}_e \le n$ is not.
\end{proposition}
\begin{proof}
Extend the interpretation with $\mathsf{dse}_e = 0$.
\end{proof}

That is, a window lowered without the flag is satisfied by a patient for whom
the event never occurred --- and ``the test was never done'' is the commonest
real violation of a monitoring recommendation.\\

\begin{proposition}[A free baseline cannot certify]
\label{prop:baseline}
Let the consequent be a reduction of $v$ by a fraction $p$ measured from
$v@e$, with $v@e$ unpinned, $\mathcal{D} \models v@e > 0$, and $P$ pinning $v = c$ for some
$c>0$. If the antecedent is satisfiable then
$\mathcal{D}\wedge P\wedge\neg\Phi_R$ is satisfiable, so
$\mathrm{verdict}(R,P) \neq \textsc{certified-compliant}$.
\end{proposition}
\begin{proof}
Choose $v@e$ in $\bigl(0,\ c/(1-p)\bigr)$, which is non-empty as $c>0$ and
$0<p<1$. Then $(1-p)\cdot v@e < c = v$, so the reduction constraint fails and
$\neg\Phi_R$ holds.
\end{proof}

A relative-change rule whose baseline no query supplies therefore behaves
exactly like a rule with an undefined predicate, which is the intended
behaviour: the lowering emits the baseline as a declared null, and the rule
answers \textsc{may-violate} until a site binds it.

\subsection{Why non-degeneracy is checked before a rule is admitted}

\begin{proposition}[Vacuity is invisible at the verdict]
\label{prop:vacuous}
If $\mathcal{D} \models \Phi_R$ then $\mathrm{verdict}(R,P) =
\textsc{certified-compliant}$ for every $P$ with $\mathcal{D}\wedge P$
satisfiable.
\end{proposition}
\begin{proof}
$\mathcal{D}\wedge P\wedge\neg\Phi_R$ is unsatisfiable.
\end{proof}

A rule that forbids nothing is therefore indistinguishable, from its output
alone, from a rule that every patient happens to satisfy. It does not crash and
it passes every syntactic gate. This is why the pipeline runs the satisfiability
of $\neg\Phi_R$ as an admission check rather than reporting it as a warning, and
why duplicate names denoting one concept are merged first: two symbols for one
concept can make $(A \wedge B) \rightarrow B$ out of a rule that was meant to
constrain something.
